\documentclass[10pt,a4paper]{amsart}
\usepackage[margin=19mm]{geometry}
\usepackage{amsmath,amssymb,mathtools}
\usepackage[scr=rsfs,cal=euler]{mathalfa}
\usepackage{xfrac}
\usepackage[unicode,hidelinks,bookmarks=false]{hyperref}
\newcommand{\ie}{i.e.\ }
\newcommand{\cf}{cf.\ }
\newcommand{\resp}{respectively\ }
\newcommand{\Subsection}{\subsection}
\providecommand{\nobreakdash}{\nobreak\hbox{-}}
\setlength{\emergencystretch}{2em}
\multlinegap0pt
\usepackage{bm}
\usepackage{bbm}
\usepackage{dsfont}
\usepackage{xspace}
\usepackage{cancel}
\usepackage{mathtools}
\usepackage{amsthm}
\makeatletter
\@ifundefined{theorem}{\newtheorem{theorem}{Theorem}}{}
\@ifundefined{lemma}{\newtheorem{lemma}[theorem]{Lemma}}{}
\makeatother
\makeatletter
\def\cen{\circledast}
\newcommand{\ms}[1]{\mathscr{#1}}
\newcommand{\ova}[1]{\overrightarrow{#1}}
\expandafter\ifx\csname proofclaim\endcsname\relax
\newenvironment{proofclaim}[1][Proof of claim]{\begin{proof}[#1]\renewcommand*{\qedsymbol}{\(\blacksquare\)}}{\end{proof}}
\fi
\makeatother
\title[Rainbow subgraphs of star-coloured graphs]{Rainbow subgraphs of star-coloured graphs}
\author{Allan Lo, Klas Markstr\"om, Dhruv Mubayi, Katherine Staden, Maya Stein, Lea Weber}
\date{Source: arXiv:2511.12505v1 (CC BY 4.0)}
\begin{document}
\maketitle

\noindent\emph{Source and licence.} Rainbow subgraphs of star-coloured graphs. Version arXiv:2511.12505v1 (CC BY 4.0), distributed under the Creative Commons Attribution 4.0 International licence, as recorded in the arXiv licence metadata for this exact version and shown on the abstract page https://arxiv.org/abs/2511.12505v1. Full rights notice: licenses/arxiv-2511.12505-CC-BY-4.0.txt.\par\smallskip

\noindent\textit{Editorial scope.} Counted unit: the Lemma whose source label is \texttt{lm:hamcycle} in the article (source0.tex lines 845--848, source label \texttt{lm:hamcycle}), delivered together with the article's own complete original proof of it (source0.tex lines 851--900, one \texttt{proof} environment of 50 lines). No mathematics is written, repaired or re-derived: statement and proof are the article's own text line for line. Required context retained uncounted: lm:ext (lines 835--838). Every definition, display and label of those lines is retained unchanged. Omitted from this package and not redistributed: 1--834, 839--844, 849--850, 901--2236. Nothing inside a retained excerpt is omitted or altered: every retained body is delivered exactly as the article's lines are written, so no omission receipt of any kind accompanies this package. Citations: the retained excerpts carry no citation command at all, so no bibliography entry of the article is transcribed and no reference is redistributed. Reference adaptation: none was needed and none was made; every reference inside the delivered unit resolves inside it, so no unresolved reference or citation is produced. Printed slips kept literal: no printed word, symbol, hypothesis or number of the counted statement or of its proof was altered. Layout and class: the article's own class is \texttt{article} with packages including multicol, amssymb, amsmath, amsthm, bm, hyperref, mathrsfs, verbatim, caption, cleveref, graphicx, float, color, mathtools; the wrapper is the lane's pinned \texttt{amsart} editorial wrapper at 10pt on A4, which restates the theorem-like environments the retained text needs and adds the article's own macro definitions that the retained text needs (\texttt{\textbackslash cen}, \texttt{\textbackslash ms}, \texttt{\textbackslash ova}), with extra wrapper packages (bm, bbm, dsfont, xspace, cancel, mathtools). The delivered numbering is the unit's own. Assets: no retained span contains a figure environment, a graphics-inclusion command, a PDF-page inclusion, a table or a TikZ picture; no image file is copied into this package, the package declares no asset dependency and no image file is redistributed with it. Licence: version arXiv:2511.12505v1 of this article is distributed under the Creative Commons Attribution 4.0 International licence, as recorded in the arXiv licence metadata for this exact version and shown on the abstract page https://arxiv.org/abs/2511.12505v1. A full rights notice is delivered at licenses/arxiv-2511.12505-CC-BY-4.0.txt. Characters: every retained excerpt is pure ASCII, so the delivered engine, XeLaTeX with its default Latin Modern text font, reports no missing character.
\begin{lemma}\label{lm:ext}
Let $n \geq 3$ be an integer, let $\ms{G}$ be a star-colouring of $K_n$ and let $v\in V(K_n)$.
	If $v$ is the centre of at least two stars and $\ms{G}-v$ has a rainbow $C_{n-1}$, then $\ms{G}$ has a rainbow $C_n$. 
\end{lemma}

\begin{lemma}\label{lm:hamcycle}
Let $n \geq 3$ be an integer and let $\ms{G}$ be a star-coloured $K_n$.
	If every vertex is the centre of at least two stars, then $\ms{G}$ has a rainbow $C_n$. 
\end{lemma}

\begin{proof}
    We proceed by induction on $n$.
	The statement is true for $n=3$, and we let $n \geq 4$ and assume that it is true for graphs with at most $n-1$ vertices.
    Recall that we write $\cen(x)$ for the number of stars centred at a vertex $x$, and $\cen_1(x)$ for the number of single-edge stars centred at~$x$.

    If there is a vertex $x$ so that $\ms{G}-x$ has at least two stars centred at each vertex, then by induction, $\ms{G}-x$ has a rainbow $C_{n-1}$. By Lemma \ref{lm:ext} we are done.   
    Thus we may assume that for every vertex $v$, there is a vertex $u$ such that $uv$ is a single-edge star and $1 \leq \cen_1(u) \leq \cen(u)=2$.

    Let $A$ be the set of vertices $u$ with $1 \leq \cen_1(u) \leq \cen(u)=2$.
    We say that a triangle is \emph{bad} if its edges are single-edge stars and it has at least two vertices in $A$.
    We claim that there is at most one bad triangle. 
    Suppose, for contradiction, that $v_1u_1u'_1$, $v_2u_2u'_2$ are two bad triangles with $u_1,u_1',u_2,u_2' \in A$. 
    By the definition of~$A$, we have $\cen_1( u ) = \cen( u ) = 2$ for $u \in \{ u_1, u_1',u_2,u_2'\}$.
    Thus, for any $i \in \{ 1,2\}$ and all vertices $w \notin \{ v_i,u_i,u'_i\}$, the edge $w u_i$ is in the star centred at~$w$.
    Hence $\{u_1,u'_1\} = \{u_2,u_2'\}$ and so $ v_1 \ne v_2$ implying that $\cen_1(u_1) \ge |\{v_1,v_2, u_1'\}| = 3$, a contradiction. 
    Let $B$ be vertices of the unique bad triangle if it exists, otherwise $B := \emptyset$.


    Form an auxiliary digraph~$\ova{J}$ with vertex set~$V(\ms{G})$ by adding the arc $\ova{uv}$ whenever $u$ is the centre of exactly one star in $\ms{G}-v$; that is, when $uv$ is a single-edge star and $u \in A$.
    We have $d^-_{\ova{J}}(v) = d^-_{\ova{J}}(v,A) \geq 1$ for all $v \in V(\ms{G})$ and $1 \le d^+_{\ova{J}}(u) \le 2 $ for all $u \in A$. 
    
    Suppose that $d^-_{\ova{J}}(v) \geq 2$ for all $v \notin B$.
    Let $L$ be the set of arcs of $\ova{J}$ whose endpoint is not in $B$. Note that $L$ contains no arc whose startpoint is in $A \cap B$. 
    Thus
    \begin{align*}
        2(n - |B|) \le \sum_{v \notin B} d^-_{\ova{J}}(v) \le 
        |L| \le \sum_{u \in A \setminus B} d^+_{\ova{J}}(u) \le 2 |A \setminus B|.
    \end{align*}
    Hence, we have $V(\ms{G}) \setminus B = A \setminus B$ and, for all $v \in V(\ms{G}) \setminus B$, we have $d^+_{\ova{J}}(v) = d^-_{\ova{J}}(v) = 2$. 
    Note that if $\ova{u w} \in E( \ova{J} )$ with $u,w \in A$, then $\ova{w u} \in E( \ova{J})$.
    Together with $B$, we deduce that $|A| \ge n- |B|$,  that $\ms{G}$ contains a $2$-factor~$F$, whose edges are all single-edge stars, and that $\cen_1( u ) = \cen( u ) = 2$ for all $u \in A$.
    Since $n \ge 4$ and $|B| \in \{0,3\}$, there exist $u,v \in A$ such that $uv \notin E(F)$. 
    Without loss of generality, $uv$ is in the star centred at $u$, so $\cen( u ) \ge |N_F(u) \cup \{v\}| = 3$ contradicting the fact that $u \in A$. 
    
        
    Therefore, there is a vertex $v \notin B$ with $d^-_{\ova{J}}(v)=1$.
    Let $u \in A$ be such that $\ova{uv} \in E(\ova{J})$.
    Let $z$ be a leaf of another star centred at~$v$ such that $uz$ is not a single-edge star.
    Indeed, we can find such a $z$ because by the hypothesis of the lemma, $v$ is the centre of at least one other star.
    Since $u \in A$, there is at most one $x \neq v$ such that $ux$ is a single-edge star, so we are done unless $\cen_1(v) = \cen(v) = 2$.
     Then $v \in A$ and $vu,vx,ux$ are single-edge stars, so $B = \{x,v,u\}$, a contradiction to $v \notin B$.

    We will form $\ms{G}'$, a new star-coloured $K_{n-1}$, by deleting $v$ and colouring $uz$ with a new colour. 
    If $x$ is a vertex such that the number of stars in~$\ms{G}'$ for which it is a centre is at most one, then $\ova{xv} \in E(\ova{J})$ and so $x=u$.
    But $u$ is in two centres since $uz$ is not a single-edge star in~$\ms{G}$.
    Thus $\ms{G}'$ has at least two stars centred at each vertex and by induction has a rainbow $C_{n-1}$.
    If this $C_{n-1}$ does not contain~$uz$, the rainbow~$C_{n-1}$ is a subgraph of $\ms{G}-v$ and we can apply Lemma~\ref{lm:ext}.
    If it does contain $uz$ we replace it with $uv$ and $vz$, which are from distinct stars with centre~$v$ and so do not appear in the rainbow $C_{n-1}$. 
    Thus we obtain a rainbow~$C_n$.
\end{proof}

\end{document}
